\section{Classical machinery and the model-specific boundary}
\label{sec:classical}

The general ingredients of SD-C24 are classical.  Determinant and primitive
cycle expansions for shifts go back to \citet{BowenLanford1970}; the broader
periodic-orbit zeta framework is developed systematically by
\citet{ParryPollicott1990}.  Twisted Perron--Frobenius operators and dynamical
\(L\)-functions appear in \citet{AdachiSunada1987}.  Voltage assignments
organize graph coverings in \citet{GrossTucker1977}, and graph-cover zeta
functions with multivariable weights are treated by
\citet{StarkTerras1996}.  We use this language to locate the construction,
not to claim any of these mechanisms as new.

Infinite covers require a trace distinction.  The \(L^2\)-zeta framework of
\citet{Clair2009} and the amenable group-extension approximation studied by
\citet{Sharp2020} are close conceptual precedents for traces per fundamental
domain.  In SD-C24 the regular deck group is the countable abelian group
\(\Q_{>0}^{\times}\), and the decisive calculation is unusually stark: no
base cycle has identity holonomy.  The resulting semifinite determinant is
locally one, while the ordinary lifted operator remains noncompact.

The cohomological vocabulary is also standard.  Livšic theory relates
periodic data and cocycle coboundaries under dynamical regularity hypotheses;
see, for example, the matrix-cocycle theorem of \citet{Kalinin2011}.  We do
not invoke such a theorem on the present countable graph.  The exact identity

\[
 q(n,d)=\frac nd\left(1+\frac1n\right)
\]

is proved directly and contains its own gauge boundary.

For operator determinants we rely only on standard trace-ideal facts after
proving trace class for the concrete matrix; \citet{Simon1977} is the
reference point.  Recent generalized weighted zeta formulas on finite
digraphs, such as \citet{IshikawaMorita2026}, show that weighted Euler,
Hashimoto, and Ihara expressions remain active adjacent machinery.  They do
not supply the successor--divisor cofactor classification or the Schatten
phase diagram below.

\begin{table}[t]
\centering
\small
\caption{Nearest classical mechanisms and the boundary of the present
claims.  Candidate-specific statements are proved internally.}
\label{tab:literature}
\begin{tabularx}{\textwidth}{L{0.26\textwidth}YY}
\toprule
Source family & Shared mechanism & Difference here \\
\midrule
Bowen--Lanford; Parry--Pollicott
& trace, primitive, and repetition expansions
& concrete countable graph, cofactor classes, and sharp nuclear domain \\
Adachi--Sunada
& character-twisted dynamical operators
& source-derived multiplicative cofactor and exact class census \\
Gross--Tucker; Stark--Terras
& voltages, covers, and cycle weights
& infinite successor--divisor grammar and endpoint operator \\
Clair; Sharp
& infinite/group extensions and normalized traces
& empty neutral periodic sector and separate ordinary noncompactness \\
Kalinin
& cocycle cohomology and periodic data
& elementary exact gauge identity; no Livšic theorem invoked \\
Simon
& trace ideals and Fredholm determinants
& candidate-specific iff threshold proved from the matrix \\
Ishikawa--Morita
& generalized weighted finite-digraph zeta
& countable graph and two-parameter Schatten theorem \\
\bottomrule
\end{tabularx}
\end{table}

The scoped search behind this paper combined the phrases
``successor--divisor graph,'' ``cofactor cocycle,'' ``multiplicative
holonomy,'' ``group extension of a countable Markov shift,'' and ``twisted
graph zeta,'' including 2024--2026 records.  We found no directly comparable
analysis that combines

\[
 Q(\gamma)=\prod(1+1/n),
 \qquad Q=2\iff\gamma=C_k,
\]

with the exact two-parameter trace-class domain proved here.  This is a
search-bounded statement, not a proof of priority.  In particular, we do not
claim invention of character twists, voltage graphs, group traces, cocycle
gauges, or Fredholm trace logarithms.

Two further boundaries matter.  First, \(L_{s,u}\) is a weighted vertex
adjacency on \(\ell^2(V)\), not a Ruelle operator on a Hölder function space.
Second, a coefficient of the connected trace logarithm records one primitive
class and its repetitions; a coefficient of the determinant itself can mix
products of distinct primitives.  All holonomy resolution below is therefore
performed on the normalized germ of \(-\log D\).
